\section{천이백만 배를 버리는 관}

\claimx{바깥 원천이 학습에 안 쓰인다는 진단의 마지막 문턱은 자료가 아니라 예산 상수다 --- 디스크의 이만 이천칠백여 행 가운데 하나만 모형에 닿는다}{남은 조각을 전부 받았을 때 모형이 보는 행 수가 늘어나면 이 주장은 떨어진다.}

바깥 원천의 사슬을 끝에서 끝까지 적는다. 디스크 $38{,}866{,}835$ 행이 정제를 지나 $670{,}118$ 문서가 되고($1.724138\%$), 개체 연결과 창 덮기를 지나 $35{,}641$ 삼중쌍이 되고($0.0917\%$), 모형이 실제로 보는 것은 겹당 $1{,}710$ 행이다($0.0044\%$). \textbf{디스크 $22{,}729.1$ 행 가운데 하나다.} 앞의 세 수는 앞 노트들의 커밋된 산출물에 있었고 넷째 수는 배관 상수에 있었는데, \textbf{그 넷을 한 줄에 나란히 적은 문서가 없었다.}

넷째 수는 자료 한계가 아니다. 그것은 학습 행 예산 $N_B$ 에 바깥 원천 비율을 곱해 반올림한 값이고, $N_B$ 는 소스 한 줄의 상수다. \textbf{남은 사백오십육 조각을 전부 받으면 셋째 수는 커지고 넷째 수는 그대로다.} 그러므로 「바깥 원천을 더 받자」는 처방은 이 상수가 서 있는 한 아무것도 안 바꾼다.

\section{요인을 하나만 흔든다 --- 도메인 할당량}

\claimx{학습 혼합을 유보 혼합에 맞추면 강한 정칙화 쪽에서 유보 순위상관이 짝 표준오차의 두 배를 넘게 오르고, 겹 씨앗 사이의 흔들림이 세 배 준다}{두 팔의 차가 어느 정칙화 값에서도 짝 표준오차의 두 배를 못 넘으면 이 주장은 떨어진다.}

두 팔은 예산도 같고 난수 차례도 같고 안쪽 학습 행은 \textbf{바이트로 같다}. 다른 것은 바깥 원천 쪽에 건 \textbf{도메인 할당량} 하나뿐이다 --- 유보의 도메인 몫을 목표로 두고 같은 순열을 같은 차례로 훑되 할당량이 찰 때까지만 받는다. 배선은 대조 팔이 앞 노트의 선택 함수와 \textbf{바이트로 같은 행}을 내고 그 예측이 앞 노트의 학습 함수와 \textbf{바이트로 같다}는 것을 보인다 --- 배관을 안 바꾸고 선택만 바꿨다는 증거다.

학습 혼합과 유보 혼합의 피어슨 상관이 $-0.1607$ 에서 $1.0$ 으로 간다. 강한 정칙화에서 유보 순위상관이 $0.263750$ 에서 $0.319593$ 으로 올라 차가 $+0.055843$ 이고 짝 표준오차 $0.024067$ 의 $2.3203$ 배다. \textbf{약한 정칙화에서는 안 움직인다}($-0.003447$, $0.1767$ 배). 두 정칙화 값 가운데 \textbf{하나}만 통과한다 --- 그대로 적는다. 그리고 겹 씨앗 다섯 사이의 표준편차가 $0.04449$ 에서 $0.014437$ 로 준다. \textbf{점추정이 안 움직이는 칸에서도 이것은 이득이다} --- 같은 자료로 같은 답을 더 안정하게 낸다.

\section{예산을 흔든다 --- 그리고 두 결과가 서로를 설명한다}

\claimx{성능은 예산과 함께 단조로 오르고 혼합 맞춤의 이득은 예산과 함께 사라진다 --- 혼합 맞춤은 예산의 대체재이고 그 대체는 자료가 모자란 쪽에서만 산다}{예산을 열여섯 배로 늘려도 성능이 안 오르거나, 혼합 맞춤의 이득이 큰 예산에서도 그대로면 이 주장은 떨어진다.}

예산을 $450$ 에서 $28{,}800$ 까지 여섯 배로 쓸었다. 대조 팔의 유보 순위상관이 강한 정칙화에서 $0.119153$ 에서 $0.365682$ 로 \textbf{단조로 오른다} --- 현행 예산은 그 곡선의 \textbf{중간}이다. 같은 곡선 위에서 혼합 맞춤의 이득은 $+0.179927$ 에서 $-0.000864$ 로 \textbf{사라진다}. 현행 예산의 사분의 일에서 층화 팔이 내는 값이 현행 예산의 대조 팔을 이미 넘는다.

사라지는 까닭이 자료에 있다. 유보가 요구하는 도메인 몫을 바깥 원천이 어느 지점부터 \textbf{못 댄다}. 처음 묶이는 도메인은 그 원천에 $1{,}582$ 행뿐인 도메인이고 예산 $3{,}600$ 부터다. 예산 $28{,}800$ 에서는 \textbf{열 도메인 가운데 여덟}이 공급에 묶이고, 그때 층화 팔은 대조 팔과 같은 팔이 된다. \textbf{그러므로 두 처방은 경쟁이 아니라 순서다} --- 자료가 모자라면 혼합을 맞추고, 자료가 넉넉하면 예산을 늘린다.

\section{크기 관문을 원리상 못 넘는 칸을 적는다}

\claimx{약한 정칙화 칸에서는 예산을 열여섯 배로 늘려 사는 전부가 결정 게이트의 크기보다 작다 --- 그 칸의 결론은 「효과 없음」이 아니라 「이 설계로는 못 잰다」이다}{그 칸의 오라클 천장이 결정 게이트 크기보다 크면 이 주장은 떨어진다.}

크기 칸을 측정 전에 등록했다. 채택 크기는 \textbf{못 정했다}고 적었다 --- 「혼합을 맞추면 얼마나 올라야 쓸모가 있나」를 정할 근거가 저장소에 없다. 대신 오라클 천장의 \textbf{형식}을 등록했다: 예산을 열여섯 배로 늘려 사는 전부다. 강한 정칙화에서 그 천장은 $0.101932$ 이고 결정 게이트 크기 $0.043972$ 의 두 배 넘게 크다 --- 비가 $0.4314$ 다. 약한 정칙화에서는 천장이 $0.017942$ 이고 게이트 크기가 $0.037700$ 이라 \textbf{비가 $2.1012$} 다. \textbf{그 칸에서는 어떤 혼합 팔도 관문을 원리상 못 넘는다.} 그 칸의 음수 값을 「혼합 맞춤은 효과가 없다」로 읽으면 그것은 \textbf{어떤 자료로도 반증할 수 없는 문장}이 된다.

\section{자를 고정하고, 자를 고치는 일을 잠근다}

\claimx{앞 세 노트가 각각 자기가 쓴 규칙으로 정본 자를 갈았고 그 마지막 선택은 자기 뽑기 잡음보다 작은 차이로 갈렸다 --- 그래서 규칙을 안 고치고 입력에서 잡음만 뺀다}{그 선택을 가른 차이가 그 양의 뽑기 잡음보다 크면 이 주장은 떨어진다.}

앞 노트의 정본 자는 도메인 안 순열 귀무의 표준편차 제곱의 역수로 가중한 자였고, 그 자와 \textbf{닫힌 꼴} 판을 가른 것은 동률 규칙이었다. 두 자의 가장 큰 도메인 몫 차는 $0.003922$ 인데 \textbf{그 몫의 뽑기 잡음이 $0.005821$} 이다 --- 비가 $0.674$ 다. \textbf{이긴 차가 자기 잡음보다 작다.} 게재된 비용도 잡음 판의 수였다.

닫힌 꼴은 항등식이므로 두 자는 \textbf{같은 자}이고, 하나는 그 값을 이천 번 뽑아 추정한 판, 하나는 값 자체다. 그래서 이 노트는 선택 규칙의 문언을 \textbf{한 글자도} 안 바꾸고 규칙에 넣는 입력에서 몬테카를로 오차만 뺀다. 그러면 씨앗 의존이 사라지고 동률 무리가 소멸하며 검정력은 $0.003\%$ 만 달라진다. \textbf{그리고 다음 노트부터 무는 조항을 신설했다} --- 한 사이클은 자기가 적용할 선택 규칙을 자기가 개정할 수 없고, 규칙은 적용하는 사이클보다 최소 한 사이클 앞서 등록한다. \textbf{자를 영구히 고정할 증거가 무엇인지는 「못 정했다」로 적었다} --- 그러므로 이 이동은 잠정이다.
